\section{Canonical source register}
\label{sec:source-register}

This register fixes the bibliographic identity, version policy, and evidential
role of every source used to define Siegel's canonical research corpus.  A
``lineage'' below is an arXiv identifier, not a claim that every version under
that identifier is the same mathematical text.  Likewise, an author-site
``last modified'' date is provenance metadata, not a rule of mathematical
precedence.

\subsection{Authority order}

Precedence is decided claim by claim, in the following order.

\begin{enumerate}[label=\textup{(\arabic*)}]
\item An explicit authorial correction or forward convention controls the
      object it names.  Thus 2601.17030v3 controls the definitions of Hydra
      maps and numens: it warns that earlier definitions, including the
      dissertation's, may be inconsistent and declares its definitions the
      standard going forward.  For frames, F-series, and rising-continuity,
      2307.00436v1 explicitly takes precedence where it conflicts with the
      dissertation.
\item Subject to (1), the latest substantive, non-withdrawn arXiv version
      controls the theorem stated by that logical work.  A metadata-only
      withdrawal revision contains no replacement mathematics.
\item The dissertation supplies the largest continuous exposition and all
      material not replaced by a later specific source.  Its three public
      witnesses are collated rather than silently identified.
\item A journal-layout PDF is a same-work witness.  It is not presumed to be
      textually identical to the arXiv source, and typography alone does not
      make it a later mathematical correction.
\item An author-site post controls its own expository wording and any material
      unique to that post.  It does not supersede an arXiv theorem merely
      because the page reports a later modification date.
\item A materially distinct earlier version remains part of the critical
      corpus.  A withdrawn work remains part of the historical corpus, but it
      cannot certify a current theorem.
\end{enumerate}

This order fixes source selection; it does not certify truth.  A controlling
witness may still contain a typographical error, an omitted hypothesis, or a
proof sketch, each of which is recorded without alteration.

\subsection{The ten arXiv lineages}

Submission dates below are UTC dates in the archived arXiv metadata.  The
similarity figures are derived TeX-line routing metrics, not altered source
witnesses.  The reproducible procedure is the following.  For the first TeX
file in each version profile, 
\path{scripts/profile_version_deltas.py} reads UTF-8 text with replacement of
decoding errors, splits it into lines, deletes the suffix beginning at an
unescaped percent sign, replaces each nonempty whitespace run by one space,
trims the result, and discards empty lines.  If the resulting line sequences
have lengths $N_1,N_2$ and the matching blocks returned by Python's
\texttt{difflib.SequenceMatcher} with \texttt{autojunk=False} contain $M$
lines, the reported value is
\[
 \frac{2M}{N_1+N_2}.
\]
The inputs are
\path{public_records/maxwell_siegel_arxiv_2026-08-16.json} and
\path{dossiers/_structural/LATEX_PROFILE_INDEX.jsonl}; exact compared paths,
source SHA-256 hashes, line counts, edit counts, and full-precision values are
in \path{dossiers/_structural/VERSION_DELTAS.jsonl}, whose SHA-256 is
\nolinkurl{40520bb04b54100bbc3a640f752c9fba0fbba4378545b3ad8b01dbb459d22647}.
This derived diagnostic only orders collation work.  It neither changes a
source nor proves semantic equivalence, semantic difference, or permission to
merge witnesses.

\begingroup\small
\setlength{\tabcolsep}{3pt}
\begin{longtable}{@{}p{0.12\textwidth}p{0.31\textwidth}p{0.25\textwidth}p{0.25\textwidth}@{}}
\toprule
Lineage & Work & Witness history & Canonical treatment\\
\midrule
\endfirsthead
\toprule
Lineage & Work & Witness history & Canonical treatment\\
\midrule
\endhead
1909.09733 & \emph{Conservation of Singularities in Functional Equations Associated to Collatz-Type Dynamical Systems; or, Dreamcatchers for Hydra Maps}
& v1, 20 Sep. 2019; v2, 4 Oct. 2019; v3, 19 Oct. 2019.  Similarities: v1--v2, $0.6041$; v2--v3, $0.8736$.
& Use v3; retain v1 and v2 as substantial same-title witnesses.  This is the earlier unit-disk/profinite dreamcatcher branch, not a chapter-version of the dissertation's central $(p,q)$-adic branch.\\

2002.12762 & \emph{Fourier Analysis of the Parity-Vector Parameterization of the Generalized Collatz $px+1$ maps}
& v1, 24 Feb. 2020, substantive paper; v2, 14 Dec. 2023, metadata-only withdrawal.
& Withdrawn as obsolete and under the wrong subject classification.  Preserve v1 only as a parity/Fourier precursor and historical witness.\\

2007.15936 & \emph{The Collatz Conjecture \& Non-Archimedean Spectral Theory: Part I---Arithmetic Dynamical Systems and Non-Archimedean Value Distribution Theory}
& v1, 19 Jul. 2020, titled \emph{Syracuse Random Variables and the Periodic Points of Collatz-type maps}; v2, 28 Mar. 2023; v3, 24 Apr. 2023.  Similarities: v1--v2, $0.0248$; v2--v3, $0.9790$.
& Use v3 for the shortened-$qx+1$ numen and Correspondence Principle.  Treat v1 as a materially distinct predecessor, not a minor draft.  Collate the journal-layout PDF and Web installments I--II.\\

2111.07882 & \emph{Functional Equations Associated to Collatz-Type Maps on Integer Rings of Algebraic Number Fields}
& v1, 10 Nov. 2021, substantive paper; v2, 16 Apr. 2022, metadata-only withdrawal.
& Withdrawn for an error in Section 3 concerning a quotient of a ring of algebraic integers modulo an ideal.  Historical only; do not use it to certify the number-field extension.\\

2111.07883 & \emph{The Collatz Conjecture \& Non-Archimedean Spectral Theory: Part I.5--How To Write The (Weak) Collatz Conjecture As A Contour Integral}
& v1, 10 Nov. 2021, titled \emph{A $p$-Adic Characterization of the Periodic Points of a Class of Collatz-Type Maps on the Integers}; v2, 15 Dec. 2023.  Similarity $0.0551$.
& Use v2 for the Dirichlet/Mellin--Perron branch.  Preserve v1 separately as a substantially different foundational witness to the Correspondence Principle.\\

2208.11082 & \emph{The Collatz Conjecture \& Non-Archimedean Spectral Theory: Part II---$(p,q)$-Adic Fourier Analysis and Wiener's Tauberian Theorem}
& v1, 18 Aug. 2022, titled \emph{Wiener Tauberian Theorems and the Value Distributions of $(p,q)$-adic Functions}; v2, 20 Jun. 2025.  Similarity $0.2870$.
& Use v2; preserve v1 as a materially distinct shorter witness.  Collate Web installment III.\\

2307.00436 & \emph{Infinite Series Whose Topology of Convergence Varies From Point to Point}
& v1, 1 Jul. 2023.
& Use v1 for frames, F-series, rising-continuity, and variable topology.  It expressly takes precedence over conflicting dissertation terminology.\\

2412.02902 & \emph{$(p,q)$-adic Analysis and the Collatz Conjecture}
& v1, 3 Dec. 2024; arXiv witness of the 2022 University of Southern California dissertation.
& Use the TeX as the machine-readable backbone, collated with the August and corrected October 2022 site PDFs.  Later specific sources override it only within their declared scope.\\

2507.13358 & \emph{Algebras of $p$-Adic Distributions Induced by Pointwise Products of F-Series}
& v1, 1 Jul. 2025.
& Use v1.  The March 2025 Kansas talk is a precursor and expository witness, not a replacement source.\\

2601.17030 & \emph{The Hydra Map and Numen Formalisms for Collatz-Type Problems}
& v1, 19 Jan. 2026; v2, 29 Jan. 2026; v3, 12 Feb. 2026.  Similarities: v1--v2, $0.9691$; v2--v3, $0.9935$.
& Use v3 as the present definition authority and global-field manual.  The v3 metadata states that minor errors in Proposition 7.6 were rectified.\\
\bottomrule
\end{longtable}
\endgroup

The changed titles and direct collation of 2007.15936v1, 2111.07883v1, and
2208.11082v1 require separate witness records.  The routing values above only
prioritize that collation: neither a shared arXiv identifier nor a line metric
is a semantic merge instruction.

\subsection{The Web installments and the paper-series numbering}

The author site's four-part numbering is not the paper-series numbering.  To
prevent false citations, this edition always prefixes the site objects with
\emph{Web installment}.

\begingroup\small
\setlength{\tabcolsep}{3pt}
\begin{longtable}{@{}p{0.12\textwidth}p{0.30\textwidth}p{0.20\textwidth}p{0.31\textwidth}@{}}
\toprule
Site label & Exact title & Page dates & Relation and status\\
\midrule
\endfirsthead
\toprule
Site label & Exact title & Page dates & Relation and status\\
\midrule
\endhead
Web I & \emph{The Collatz Conjecture \& Non-Archimedean Spectral Theory--Part I--The Numen of a Hydra Map}
& Published 2022-08-09 21:45:22 UTC; page reports modification 2024-07-11 21:24:00 UTC.
& Expository witness for the construction half of paper Part I, arXiv:2007.15936v3.  It is not an additional arXiv Part I.\\

Web II & \emph{The Collatz Conjecture \& Non-Archimedean Spectral Theory--Part II--The Correspondence Principle}
& Published 2022-08-10 20:55:28 UTC; page reports modification 2022-10-11 19:55:07 UTC.
& Expository witness for the Correspondence-Principle half of \emph{paper Part I}.  Its final Dirichlet/Perron aside is a precursor to paper Part I.5, arXiv:2111.07883v2.  It is not paper Part II.\\

Web III & \emph{The Collatz Conjecture \& Non-Archimedean Spectral Theory--Part III--$(p,q)$-adic Fourier Analysis \& Tauberian Spectral Theory}
& Published 2022-08-12 22:51:30 UTC; archived page reports modification 2025-12-26 22:12:20 UTC.
& Evolving site exposition overlapping \emph{paper Part II}, arXiv:2208.11082v2.  Unique material may be used after theorem-level collation; its modification date alone confers no authority.\\

Web IV & \emph{The Collatz Conjecture \& Non-Archimedean Spectral Theory--Part IV--A Fourier Analysis of the Numen of the Shortened $qx+1$ map}
& Published 2022-08-15 19:14:38 UTC; archived page reports modification 2025-12-26 22:13:00 UTC.
& Site-only numen/Fourier application, chiefly parallel to the dissertation's one-dimensional Fourier material.  There is no corresponding arXiv paper Part IV in this corpus.\\
\bottomrule
\end{longtable}
\endgroup

The archived Web I and Web II objects are current snapshots acquired on
23 August 2026.  WordPress's public endpoint did not expose their revision
histories without authentication.  Consequently their published and modified
timestamps do not identify the contents of every intervening revision.

\subsection{Dissertation and supplemental author-site witnesses}

The dissertation has three distinct public witnesses:

\begin{enumerate}[label=\textup{(\alph*)}]
\item the August 2022 author-site PDF, 476 pages;
\item the October 2022 author-site PDF explicitly labelled corrected,
      478 pages; and
\item arXiv:2412.02902v1, submitted 3 December 2024, 467 pages, with an
      extractable TeX source.
\end{enumerate}

Their normalized texts and pagination differ.  The arXiv TeX is therefore an
editable base, not evidence that the explicitly corrected October PDF is a
byte-for-byte or proposition-for-proposition predecessor.  Any difference
used in the edition must be collated at the relevant passage.

The remaining author-site witnesses have the following roles.

\begin{description}[style=nextline]
\item[Part I journal-layout PDF.]
Same-work layout witness for 2007.15936; use for publication pagination and
typographic variants only after comparison with v3.

\item[Infinite-series journal-layout PDF.]
Same-work layout witness for 2307.00436; it does not silently replace the
available TeX.

\item[\emph{Arithmetic Properties of F-series} talk notes.]
March 2025 Emporia State University talk; supplementary precursor to
2507.13358.

\item[Episode 1 script.]
June 2024 companion exposition for the dissertation and the $(p,q)$-adic
programme; useful orientation, not a primary theorem source.

\item[$p$-adic $L$-functions seminar and Fourier-series notes.]
Siegel-authored background with high contextual relevance.  These are
supplementary notes, not parts of the Collatz research lineage.
\end{description}

\subsection{Corpus boundary}

The core edition includes the ten arXiv lineages, all materially distinct
versions identified above, the four Web installments, the three dissertation
witnesses, the two same-work layout witnesses, and the research-specific talk
and companion script.  The $p$-adic $L$-function seminar and Fourier-series
notes are Siegel-authored apparatus, not core research works.  Authorship and
core status are separate questions.  Siegel's generic
course notes in introductory real analysis, linear algebra, vector calculus,
complex variables, partial and ordinary differential equations, and
probability remain Siegel-authored works, but they are outside the Collatz,
Hydra, numen, dreamcatcher, and F-series research corpus.

External supporting literature is not part of Siegel's corpus, but it is used
with exact citations to establish prior art, standard theorems, and comparison
maps.

\subsection{External prior-art comparison ledger}

The following sources are outside Siegel's corpus.  They are present because
an exact critical edition must distinguish Siegel's added structure from an
older identical special case, and must distinguish a shared value formula
from an identity of maps.  Each row records the strongest comparison proved
in this edition and the boundary beyond which the citation may not be used.

\begingroup\small
\setlength{\tabcolsep}{3pt}
\begin{longtable}{@{}p{0.19\textwidth}p{0.23\textwidth}p{0.28\textwidth}p{0.23\textwidth}@{}}
\toprule
Object & Primary source & Exact proved overlap & Nonidentity boundary\\
\midrule
\endfirsthead
\toprule
Object & Primary source & Exact proved overlap & Nonidentity boundary\\
\midrule
\endhead
Scalar integer Hydra at displayed modulus $D\in\Z_{\geq2}$ &
Matthews, Sec.~2 \cite{MatthewsGeneralized} & Admissibility is equivalent to
unique integers $m_i\ne0$ and $\rho_i\equiv im_i\pmod D$ with
$H_i(x)=(m_ix-\rho_i)/D$; see
\cref{prop:scalar-hydra-matthews}.  Exact-integrality is equivalent to
the displayed-presentation condition
$\gcd(m_0\cdots m_{D-1},D)=1$. & Matthews attaches relatively prime type to
the least-modulus presentation.  Thus exact-integrality is equivalent to
Matthews relatively prime type only together with
$D=d_{\mathrm{Mat}}(H)$.  A nonminimal refinement may destroy
exact-integrality.  Neither comparison identifies Siegel's non-scalar or
global-field structure or adds centeredness or properness.\\

Residue-class-wise affine mapping & Kohl, Defs.~1.1 and~1.4 and Rem.~1.2
\cite{KohlRCWA} & Every scalar integer Hydra in the preceding row is rcwa,
and $d_{\mathrm{Mat}}(H)=\operatorname{Mod}(H)\mid D$; under
exact-integrality both least moduli equal $D$; see
\cref{prop:scalar-hydra-rcwa}. & The group
$\operatorname{RCWA}(\Z)$ contains only bijective rcwa maps.  The shortened
Collatz map is rcwa but is not injective; a Hydra is therefore not
automatically an rcwa group element.\\

Formal affine word and $q=3$ return candidate & B\"ohm--Sontacchi,
Prop.~1; printed Props.~4--5; unnumbered paragraph after Prop.~5 on p.~264
\cite{BohmSontacchi} & Proposition~1 specializes to the valid formal
composition in \cref{eq:chronological-affine-return}; at $q=3$, its fixed
point is the rational $R_{3,\boldsymbol\varepsilon}$ printed as the
Proposition~5 candidate after the explicit renaming $n=L$. & Printed
Proposition~4's unrestricted identification with $u^L(x)$ is false; it may be
made only for a matching itinerary.  The unnumbered p.~264 paragraph allows
exact period properly dividing $L$.  The final strict position inequality is
reconstructed from the formal run data, not inserted into printed
Proposition~5.\\

$2$-adic parity conjugacy & Bernstein--Lagarias, Eqs.~(1.3)--(1.6),
(4.1)--(4.2), and App.~B \cite{BernsteinLagarias} & The chronological parity
encoder $Q_q$, inverse conjugacy $\Phi_q$, and every odd-$q$ classical
specialization are exact.  For odd prime $q$, the $\Phi_q$ value is
$\iota_2(R_{q,\boldsymbol\varepsilon})$ and the reversed $\chi_q$ value is
$\iota_q(R_{q,\boldsymbol\varepsilon})$ as in
\cref{thm:BL-numen-periodic-overlap}. & $Q_q$ and
$\Phi_q$ retain their inverse roles in $\Z_2$; for odd prime $q$, neither is
the cross-adic map $\chi_q:\Z_2\to\Z_q$.  Composite odd $q$ in the classical conjugacy does not
enlarge the edition's prime-$q$ numen theorem.\\

De Bruijn graph presentation & Laarhoven--de Weger, Thms.~3.1 and~5.1 and
Sec.~5.2 \cite{LaarhovenDeWeger} & Their graph isomorphism is the
chronological parity map denoted $Q_3$ here; their odd-$an+b$ analogue is an
exact classical bridge. & Their paper's symbol $\Phi$ has the orientation of
this edition's $Q_3$, not Bernstein--Lagarias's $\Phi_3$; its $p$-ary
discussion uses appropriately chosen coefficients and is not an arbitrary-
Hydra theorem.\\
\bottomrule
\end{longtable}
\endgroup

\subsection{Source-issue ledger}

\begingroup\small
\begin{longtable}{@{}p{0.10\textwidth}p{0.22\textwidth}p{0.61\textwidth}@{}}
\toprule
Code & Source & Canonicalization issue\\
\midrule
\endfirsthead
\toprule
Code & Source & Canonicalization issue\\
\midrule
\endhead
\textsc{SR-01} & Corpus index & A shared logical-work identifier records bibliographic lineage, not semantic sameness.  The reproducible routing values $0.0248$, $0.0551$, and $0.2870$ do not establish that distinction and do not authorize a merge; the separately archived versioned sources and direct collation control witness status.\\

\textsc{SR-02} & 1909.09733 & A common title is not an identity proof.  Versions v1 and v2 remain separately archived in the apparatus even though v3 is the base; the value $0.6041$ is only a reproducible collation-routing diagnostic.\\

\textsc{SR-03} & 2002.12762 & v2 is a metadata-only withdrawal declaring the paper obsolete and under the wrong subject classification.  It supplies no corrected paper; v1 is historical only.\\

\textsc{SR-04} & 2111.07882 & v2 withdraws the work for an error in Section 3 about a quotient ring.  The archived comment contains the contradictory phrase ``was not incorrectly described''; the withdrawal and existence of an error are clear, but the intended corrected assertion must not be guessed.\\

\textsc{SR-05} & 2601.17030 & The paper expressly replaces potentially inconsistent earlier Hydra definitions.  v3 also reports corrections to minor errors in Proposition 7.6.  Earlier Hydra terminology is historical unless re-derived in the v3 formalism.\\

\textsc{SR-06} & 2307.00436 & The paper expressly takes precedence over the dissertation where their frames/F-series terminology conflicts.  It is more general in some respects and less comprehensive in others; only the conflicting definitions are superseded.\\

\textsc{SR-07} & Dissertation & The August 2022, corrected October 2022, and December 2024 arXiv witnesses are textually distinct.  ``Later upload'' and ``corrected'' are recorded attributes, not a license to discard the other witnesses.\\

\textsc{SR-08} & Web numbering & Web II is part of the web sequence but belongs chiefly to arXiv \emph{paper Part I}; Web III overlaps arXiv \emph{paper Part II}.  Bare labels such as ``Part II'' are therefore prohibited in source-critical references.\\

\textsc{SR-09} & Web I & The current snapshot calls \emph{numen} Greek whereas paper Part I correctly identifies it as Latin; its displayed $q$-adic decay formula omits the subscript $q$ on the norm; and its uniqueness proof says equality on $\Z_q$ where the input domain must be $\Z_2$.  Paper Part I v3 controls these points.  Hash-bound locator: \path{author_site/wordpress_2026-08-23/part_i_numen_hydra_map_2024-07-11.txt}, SHA-256 \nolinkurl{6dfcf3ee7c4056978deae669094e08f41acbfbd7f1410b2cafb3d58293f682ae}, lines 143, 271--275, and 327--329.\\

\textsc{SR-10} & Web II & Claim 5's statement omits the conclusion $\chi_q(z/2)=0$, although its proof supplies it.  The page calls its terminal periodic-point formulation Version 3, while paper Part I v3 canonically labels its corresponding published formulation Version 4.  Its divergent result says ``belongs to a divergent trajectory''; paper Part I v3 states directly that the value is a divergent point.  Hash-bound locator: \path{author_site/wordpress_2026-08-23/part_ii_correspondence_principle_2022-10-11.txt}, SHA-256 \nolinkurl{b5d37b71f4778109fbd853f621b0562c6aeca1487141968d2cde3887c8ba46c6}, lines 283--287, 353--355, and 359--373.\\

\textsc{SR-11} & Web I--IV & Publication and modification timestamps are provenance only.  No page supersedes a formal source solely because its modification date is later, and the public snapshots do not furnish a complete revision history.\\

\textsc{SR-12} & Web III--IV & These are active, evolving expositions.  Statements unique to them require theorem-level checking before incorporation; overlap with paper Part II or the dissertation is not presumed identity.\\

\textsc{SR-13} & Layout PDFs & The Part I and infinite-series layouts are distinct files and must be collated for wording and theorem numbering.  They are not replacement TeX sources merely because they resemble journal typography.\\

\textsc{SR-14} & 2111.07883v2; Web II & In the calculation of $f_q(z^2,x)$ immediately before the printed ordinary generating-function equation, the first term is written as $z/(1-z)$ instead of $z^2/(1-z^2)$.  The resulting Dirichlet functional equation is false for $x\ne0$, and its $q=3$ residue formulas contain spurious $x$-terms.  Part I.5 and the public Web II nevertheless identify the continuation programme.  Chapter 7 recovers it by deriving the corrected equation, proving normal convergence and global strip-by-strip continuation, and recalculating the principal parts.  Denominator zeros remain candidate poles unless noncancellation is proved.\\

\textsc{SR-15} & B\"ohm--Sontacchi (1978) & Printed Proposition~4 is false
under its stated unrestricted quantifiers.  With $x=0$, $m=1$, and
$n_0=n_1=1$, it gives the right side $1/4$, while $u^3(0)=0$; the right side
is instead the formal composition $f\circ g\circ f(0)=1/4$.  The edition
therefore derives the formal affine identity from Proposition~1, invokes
$u^L$ only after proving actual-itinerary correctness, and recovers the
Proposition~5 fixed-return candidate independently.  The possibility of
exact period properly dividing $L$ is located in the unnumbered paragraph
after Proposition~5 on p.~264, not attributed generically to
Propositions~4--6.\\

\bottomrule
\end{longtable}
\endgroup

\begingroup
\small
\sloppy
\def\UrlBreaks{\do\/\do-\do\_}
\Urlmuskip=0mu plus 1mu\relax
\begin{thebibliography}{99}

\bibitem{SiegelDreamcatchers}
Maxwell C. Siegel,
\emph{Conservation of Singularities in Functional Equations Associated to
Collatz-Type Dynamical Systems; or, Dreamcatchers for Hydra Maps},
arXiv:1909.09733.  Version history: v1, 20 September 2019; v2,
4 October 2019; v3, 19 October 2019.  Canonical witness: v3.
\url{https://arxiv.org/abs/1909.09733v3}.

\bibitem{SiegelParityFourier}
Maxwell C. Siegel,
\emph{Fourier Analysis of the Parity-Vector Parameterization of the
Generalized Collatz $px+1$ maps}, arXiv:2002.12762v1, 24 February 2020;
withdrawal revision v2, 14 December 2023.
Substantive v1: \url{https://arxiv.org/abs/2002.12762v1}.
Withdrawal record v2: \url{https://arxiv.org/abs/2002.12762v2}.

\bibitem{SiegelSyracuseV1}
Maxwell C. Siegel,
\emph{Syracuse Random Variables and the Periodic Points of Collatz-type
maps}, arXiv:2007.15936v1, 19 July 2020.  Materially distinct predecessor
under the later Part I lineage.
\url{https://arxiv.org/abs/2007.15936v1}.

\bibitem{SiegelPartI}
Maxwell C. Siegel,
\emph{The Collatz Conjecture \& Non-Archimedean Spectral Theory: Part I---
Arithmetic Dynamical Systems and Non-Archimedean Value Distribution Theory},
arXiv:2007.15936v3, 24 April 2023; v2 submitted 28 March 2023.
Publication DOI: \url{https://doi.org/10.1134/S2070046624020055}.
\url{https://arxiv.org/abs/2007.15936v3}.

\bibitem{SiegelNumberRings}
M. C. Siegel,
\emph{Functional Equations Associated to Collatz-Type Maps on Integer Rings
of Algebraic Number Fields}, arXiv:2111.07882v1, 10 November 2021;
withdrawal revision v2, 16 April 2022.
Substantive v1: \url{https://arxiv.org/abs/2111.07882v1}.
Withdrawal record v2: \url{https://arxiv.org/abs/2111.07882v2}.

\bibitem{SiegelPadicCharacterizationV1}
Maxwell C. Siegel,
\emph{A $p$-Adic Characterization of the Periodic Points of a Class of
Collatz-Type Maps on the Integers}, arXiv:2111.07883v1,
10 November 2021.  Materially distinct predecessor under the later Part I.5
lineage.  \url{https://arxiv.org/abs/2111.07883v1}.

\bibitem{SiegelPartIHalf}
Maxwell C. Siegel,
\emph{The Collatz Conjecture \& Non-Archimedean Spectral Theory: Part I.5--
How To Write The (Weak) Collatz Conjecture As A Contour Integral},
arXiv:2111.07883v2, 15 December 2023.
\url{https://arxiv.org/abs/2111.07883v2}.

\bibitem{SiegelWienerV1}
Maxwell C. Siegel,
\emph{Wiener Tauberian Theorems and the Value Distributions of $(p,q)$-adic
Functions}, arXiv:2208.11082v1, 18 August 2022.  Materially distinct
predecessor under the later Part II lineage.
\url{https://arxiv.org/abs/2208.11082v1}.

\bibitem{SiegelPartII}
Maxwell C. Siegel,
\emph{The Collatz Conjecture \& Non-Archimedean Spectral Theory: Part II---
$(p,q)$-Adic Fourier Analysis and Wiener's Tauberian Theorem},
\emph{$p$-Adic Numbers, Ultrametric Analysis and Applications}
\textbf{17}(2) (2025), 187--232,
\url{https://doi.org/10.1134/S2070046625020062};
arXiv:2208.11082v2, 20 June 2025.
\url{https://arxiv.org/abs/2208.11082v2}.

\bibitem{SiegelInfiniteSeries}
Maxwell C. Siegel,
\emph{Infinite Series Whose Topology of Convergence Varies From Point to
Point}, \emph{$p$-Adic Numbers, Ultrametric Analysis and Applications}
\textbf{15}(2) (2023), 133--167,
\url{https://doi.org/10.1134/S2070046623020061};
arXiv:2307.00436v1, 1 July 2023.
\url{https://arxiv.org/abs/2307.00436v1}.

\bibitem{SiegelDissertationArxiv}
Maxwell C. Siegel,
\emph{$(p,q)$-adic Analysis and the Collatz Conjecture}, Ph.D. dissertation,
University of Southern California, 2022; arXiv:2412.02902v1,
3 December 2024.  \url{https://arxiv.org/abs/2412.02902v1}.

\bibitem{SiegelFSeriesProducts}
Maxwell C. Siegel,
\emph{Algebras of $p$-Adic Distributions Induced by Pointwise Products of
F-Series}, arXiv:2507.13358v1, 1 July 2025.
\url{https://arxiv.org/abs/2507.13358v1}.

\bibitem{SiegelHydraManual}
Maxwell C. Siegel,
\emph{The Hydra Map and Numen Formalisms for Collatz-Type Problems},
arXiv:2601.17030.  Version history: v1, 19 January 2026; v2,
29 January 2026; v3, 12 February 2026.  Canonical witness: v3.
\url{https://arxiv.org/abs/2601.17030v3}.

\bibitem{SiegelWebI}
Maxwell C. Siegel,
\emph{The Collatz Conjecture \& Non-Archimedean Spectral Theory--Part I--
The Numen of a Hydra Map}, author-site post (Web installment I in this
edition), published 2022-08-09 21:45:22 UTC; page reports modification
2024-07-11 21:24:00 UTC; snapshot acquired 23 August 2026.
\url{https://siegelmaxwellc.wordpress.com/2022/08/09/the-collatz-conjecture-non-archimedean-spectral-theory-part-i-the-numen-of-a-hydra-map/}.

\bibitem{SiegelWebII}
Maxwell C. Siegel,
\emph{The Collatz Conjecture \& Non-Archimedean Spectral Theory--Part II--
The Correspondence Principle}, author-site post (Web installment II in this
edition), published 2022-08-10 20:55:28 UTC; page reports modification
2022-10-11 19:55:07 UTC; snapshot acquired 23 August 2026.
\url{https://siegelmaxwellc.wordpress.com/2022/08/10/the-collatz-conjecture-non-archimedean-spectral-theory-part-ii-the-correspondence-principle/}.

\bibitem{FedjaMO333801}
fedja (MathOverflow user 1131),
answer 333801 to \emph{Overconcentration of Poles on the Circle of
Convergence of a Power Series with Bounded Coefficients}, accepted
MathOverflow answer, created 12 June 2019 at 03:24:50 UTC and last edited
18 June 2019 at 01:50:51 UTC, CC BY-SA 4.0.
\url{https://mathoverflow.net/a/333801}.

\bibitem{MatthewsGeneralized}
K. R. Matthews,
\emph{Generalized $3x + 1$ Mappings: Markov Chains and Ergodic Theory},
in Jeffrey C. Lagarias (ed.), \emph{The Ultimate Challenge: The $3x+1$
Problem}, American Mathematical Society, Providence, RI, 2010, pp.~79--103.
The consulted author-form witness states that it contains small additions and
deletions relative to the book chapter.

\bibitem{KohlRCWA}
Stefan Kohl,
\emph{Algorithms for a class of infinite permutation groups},
\emph{Journal of Symbolic Computation} \textbf{43}(8) (2008), 545--581.
\url{https://doi.org/10.1016/j.jsc.2007.12.001}.

\bibitem{BohmSontacchi}
Corrado B\"ohm and Giovanna Sontacchi,
\emph{On the existence of cycles of given length in integer sequences like
$x_{n+1}=x_n/2$ if $x_n$ even, and $x_{n+1}=3x_n+1$ otherwise},
\emph{Atti della Accademia Nazionale dei Lincei. Classe di Scienze Fisiche,
Matematiche e Naturali. Rendiconti}, Series~8, \textbf{64}(3) (1978),
260--264.
\url{https://www.bdim.eu/item?id=RLINA_1978_8_64_3_260_0}.

\bibitem{BernsteinLagarias}
Daniel J. Bernstein and Jeffrey C. Lagarias,
\emph{The $3x + 1$ Conjugacy Map},
\emph{Canadian Journal of Mathematics} \textbf{48}(6) (1996), 1154--1169.
\url{https://doi.org/10.4153/CJM-1996-060-x}.

\bibitem{LaarhovenDeWeger}
Thijs Laarhoven and Benne de Weger,
\emph{The Collatz conjecture and De Bruijn graphs},
arXiv:1209.3495v1, 16 September 2012.
\url{https://arxiv.org/abs/1209.3495v1}.
The consulted local PDF also displays 27 November 2024; that display is not
used as the arXiv submission date.

\bibitem{SiegelWebIII}
Maxwell C. Siegel,
\emph{The Collatz Conjecture \& Non-Archimedean Spectral Theory--Part III--
$(p,q)$-adic Fourier Analysis \& Tauberian Spectral Theory}, author-site post
(Web installment III in this edition), published 2022-08-12 22:51:30 UTC;
archived page reports modification 2025-12-26 22:12:20 UTC.
\url{https://siegelmaxwellc.wordpress.com/2022/08/12/the-collatz-conjecture-non-archimedean-spectral-theory-part-iii-a-pq-adic-fourier-analysis/}.

\bibitem{SiegelWebIV}
Maxwell C. Siegel,
\emph{The Collatz Conjecture \& Non-Archimedean Spectral Theory--Part IV--
A Fourier Analysis of the Numen of the Shortened $qx+1$ map}, author-site
post (Web installment IV in this edition), published 2022-08-15 19:14:38
UTC; archived page reports modification 2025-12-26 22:13:00 UTC.
\url{https://siegelmaxwellc.wordpress.com/2022/08/15/the-collatz-conjecture-non-archimedean-spectral-theory-part-iv-a-fourier-analysis-of-the-numen-of-the-shortened-qx-1-map/}.

\bibitem{SiegelDissertationAugust}
Maxwell C. Siegel,
\emph{$(p,q)$-adic Analysis and the Collatz Conjecture}, Ph.D. dissertation,
University of Southern California, August 2022 author-site edition.
\url{https://siegelmaxwellc.wordpress.com/wp-content/uploads/2022/08/phd_dissertation_final.pdf}.

\bibitem{SiegelDissertationOctober}
Maxwell C. Siegel,
\emph{$(p,q)$-adic Analysis and the Collatz Conjecture}, Ph.D. dissertation,
University of Southern California, corrected October 2022 author-site
edition.  \url{https://siegelmaxwellc.wordpress.com/wp-content/uploads/2022/10/phd_dissertation_final.pdf}.

\bibitem{SiegelPartILayout}
Maxwell C. Siegel,
\emph{The Collatz Conjecture \& Non-Archimedean Spectral Theory: Part I},
author-site journal-layout PDF, same-work witness for arXiv:2007.15936.
\url{https://siegelmaxwellc.wordpress.com/wp-content/uploads/2024/05/blog_post_i_numen_and_correspondence_principle.pdf}.

\bibitem{SiegelInfiniteSeriesLayout}
Maxwell C. Siegel,
\emph{Infinite Series Whose Topology of Convergence Varies From Point to
Point}, author-site journal-layout PDF, same-work witness for
arXiv:2307.00436.
\url{https://siegelmaxwellc.wordpress.com/wp-content/uploads/2023/12/infinite_series_whose_topology_of_convergence_varies_from_point_to_point_simplified.pdf}.

\bibitem{SiegelKansasTalk}
Maxwell C. Siegel,
\emph{Arithmetic Properties of F-series; or, How to $3$-adically Integrate
a $5$-adic Function}, Emporia State University talk notes, March 2025;
author-site PDF uploaded May 2025.
\url{https://siegelmaxwellc.wordpress.com/wp-content/uploads/2025/05/2025_kansas_talk-3.pdf}.

\bibitem{SiegelEpisodeOne}
Maxwell C. Siegel,
\emph{Episode 1---$(p,q)$-adic Analysis and the Collatz Conjecture---A Whole
New World: Script}, author-site companion script, June 2024.
\url{https://siegelmaxwellc.wordpress.com/wp-content/uploads/2024/06/episode-1-script.pdf}.

\bibitem{SiegelPadicLFunctions}
Maxwell C. Siegel,
\emph{$p$-adic $L$-functions}, graduate seminar notes, author-site PDF.
\url{https://siegelmaxwellc.wordpress.com/wp-content/uploads/2023/02/algebraic_number_theory_-_presentation_-_p-adic_l-functions.pdf}.

\bibitem{SiegelFourierNotes}
Maxwell C. Siegel,
\emph{Fourier Series}, course notes, author-site PDF, 2025 site copy.
\url{https://siegelmaxwellc.wordpress.com/wp-content/uploads/2025/02/math_445_souto_-_class_notes_-_chapter_11_-_fourier_analysis.pdf}.

\end{thebibliography}
\endgroup
